\documentclass[12pt]{article}
\usepackage[margin=0.8in]{geometry} 
\usepackage{setspace} 
\usepackage{float}
\usepackage{amsmath}
\usepackage{tabularx} 
\usepackage{booktabs}
\usepackage{adjustbox}
\usepackage{longtable}
\usepackage{pdflscape}
\usepackage{threeparttable} 
\usepackage{array}
\usepackage{multirow}
\usepackage{natbib}
\usepackage{amsfonts}
\usepackage{enumitem}
\usepackage{graphicx}
\setlist[itemize,1]{label=\textcolor{black}{\tiny $\bullet$}}
\setlist[itemize,2]{label=\textcolor{black}{\tiny $\bullet$}}
\setlist[itemize,3]{label=\textcolor{black}{\tiny $\bullet$}}
\usepackage{fix-cm}
\usepackage[T1]{fontenc}
\usepackage{lmodern}
\onehalfspacing
\usepackage{tcolorbox}
\tcbuselibrary{skins, breakable}

\setlength{\parindent}{0pt}
\setlength{\parskip}{6pt}

\title{\textbf{Section for Supervision}}

\author{Students: Helle and William \\
Supervisors: Arne, Mette, and Jakob}



\begin{document}

\maketitle

\section{Field-Level Profitability from the Farmer's Perspective} 
\label{sec:Farmer Profitability}

The following sections present area-weighted average DBII estimates across catchments and soil types for winter wheat and spring barley based on fields where each crop was actually cultivated in 2024. 
The reader should bear in mind the following points throughout the section: First, all estimates exclude direct CAP payments, which constitute a significant share of farm income in Denmark. Negative DBII values should therefore not be interpreted as evidence that cultivation is unprofitable for the farmer -- area-based support would shift estimates upward uniformly across all fields, however leaving the relative spatial pattern unchanged. Second, all estimates reflect 2024 conditions, which were characterised by exceptionally adverse weather -- total precipitation approximately 50 percent above the climatological normal, with yield shortfalls across all major crops compounded by lower market prices \citep{bramsen2024}. Results should therefore be interpreted as indicative of field-level profitability under 2024 conditions rather than as estimates of long-run productive value. 

Average DBII estimates are reported only for soil type--catchment combinations based 
on more than 20 fields. This threshold reflects both data confidentiality requirements 
and a judgement that estimates based on fewer observations fail to be sufficiently robust 
for the purposes of this pilot study. 



\subsubsection{Field-Level Profitability for Winter Wheat in 2024} 


\begin{table}[H]
\centering
\small
\caption{Average gross yield (DKK/ha/year) for winter wheat fields by catchment and soil type}
\label{tab:yield-vinterhvede}
\begin{threeparttable}
\begin{tabular}{lrrrrr}
\toprule
\textbf{Catchment} & \textbf{Fields} & \textbf{Yield} & \textbf{Fine} &
\textbf{Coarse} & \textbf{Loamy} \\
\midrule\addlinespace
Lister Dyb                 & 440 & 8,537  & 8,128 (117)  & 7,768 (71)   & 8,868 (252) \\
Nissum Bredning            & 169 & 8,193  & 9,125 (43)   & $-$           & 7,912 (109) \\
Nissum Fjord, Felsted Kog  & 196 & 8,153  & 8,106 (49)   & 7,786 (38)   & 8,283 (109) \\
Thisted Bredning           & 244 & 9,601  & 9,327 (50)   & 8,537 (27)   & 9,817 (167) \\
Nibe Bredning og Langerak  & 704 & 9,470  & 9,361 (229)  & 8,747 (75)   & 9,646 (400) \\
Randers Fjord, indre       & 807 & 10,080 & 10,097 (200) & 9,785 (88)   & 10,111 (519) \\
Horsens Fjord, indre       & 233 & 9,724  & 9,671 (57)   & 9,261 (32)   & 9,816 (144) \\
Odense Fjord, Seden Strand & 304 & 10,694 & 10,590 (73)  & 9,856 (40)   & 10,847 (191) \\
Isefjord, indre            & 202 & 10,928 & 11,628 (45)  & $-$           & 10,796 (139) \\
Karrebæk Fjord             & 342 & 11,098 & 11,382 (71)  & 10,264 (50)  & 11,151 (221) \\
Roskilde Fjord, indre      &  90 & 10,520 & 10,370 (23)  & $-$           & 10,759 (57) \\
\addlinespace\bottomrule
\end{tabular}
\begin{tablenotes}
\footnotesize
\item \textit{Notes:} All values are measured in DKK per hectare, and weighted by area;
Yield denotes the catchment-level area-weighted average gross yield;
Fine, Coarse, and Loamy refer to average gross yield for fields on the respective
soil type within each catchment; numbers in parentheses indicate the number of
fields; `$-$' indicates fewer than 20 fields and is therefore not reported.
\end{tablenotes}
\end{threeparttable}
\end{table}


Table~\ref{tab:yield-vinterhvede} presents average gross winter wheat yields (in DKK/ha) across the 11 catchments as of 2024. 
Substantial geographic variation is observed, revealing a clear west--–east productivity gradient. The lowest average yields are observed in western Jutland, particularly in Nissum Fjord, Felsted Kog (8,153 DKK/ha), Nissum Bredning (8,193 DKK/ha), and Lister Dyb (8,537 DKK/ha). Although Lister Dyb differs geographically from the other western catchments, as it is located in Southern Jutland and stretches from the west to the east coast, it exhibits similarly low revenue levels. 

Catchments located in northern and eastern Jutland display intermediate yield levels, including Thisted Bredning (9,601 DKK/ha), Nibe Bredning og Langerak (9,470 DKK/ha), Horsens Fjord, indre (9,724 DKK/ha), and Randers Fjord, indre (10,080 DKK/ha). The highest yields are observed in eastern Denmark and Zealand, particularly in Karrebæk Fjord (11,098 DKK/ha), Isefjord, indre (10,928 DKK/ha), Roskilde Fjord, indre (10,520 DKK/ha), and Odense Fjord, Seden Strand (10,694 DKK/ha). 

Overall, the pattern indicates  a gradual increase in winter wheat yield from the sandy and more precipitation-exposed catchments of western Denmark towards the more fertile and clay-dominated catchments of eastern Denmark and Zealand. Importantly, the west--east gradient persists even within soil categories, suggesting that soil composition alone cannot fully explain the observed spatial variation.

Differences across soil types generally follow expected agronomic patterns. Coarse soils consistently exhibit the lowest average revenues, whereas loamy soils tend to generate the highest. For example, in Nibe Bredning og Langerak, average revenues range from 8,747 DKK/ha on coarse soils to 9,646 DKK/ha on loamy soils. Similar within-catchment differences are observed across most catchments.


Descriptively, variation across catchments appears larger than variation across soil types within catchments. In several cases, coarse soils in eastern catchments outperform loamy soils in western catchments, indicating that broader geographic and climatic conditions may play an important role in shaping winter wheat production values.

A few catchments deviate somewhat from the overall geographic gradient. Karrebæk Fjord consistently exhibits the highest yield levels across all soil types and appears relatively robust given its large sample sizes. In contrast, Nissum Bredning and Nissum Fjord, Felsted Kog display persistently low yields across soil categories. Isefjord, indre shows particularly high revenues on fine soils; however, this estimate is based on a relatively limited number of observations (45) and should be interpreted with caution. Notably, Thisted Bredning and Nissum Bredning are geographically adjacent, yet Thisted Bredning exhibits considerably higher yields.

Isefjord displays particularly high yields on fine soils; however, this estimate is based on a relatively small number of observations, indicating that this result may be more sensitive to sample composition. Notably, Thisted Bredning and Nissum Bredning are geographically adjacent, yet Thisted Bredning displays considerably higher yields. 

Taken together, the relatively smooth spatial gradient may suggest, that the observed differences are not solely driven by isolated outliers or individual catchments, but instead reflect broader systematic geographic variation in winter wheat production values in 2024.

\begin{table}[H]
\centering
\small
\caption{Average total costs (DKK/ha/year) for winter wheat fields by catchment and soil type}
\label{tab:costs-vinterhvede}
\begin{threeparttable}
\begin{tabular}{lrrrrr}
\toprule
\textbf{Catchment} & \textbf{Fields} & \textbf{Costs} & \textbf{Fine} &
\textbf{Coarse} & \textbf{Loamy} \\
\midrule\addlinespace
Lister Dyb                 & 440 & 8,611  & 8,629 (117)  & 7,747 (71)   & 8,780 (252) \\
Nissum Bredning            & 169 & 10,264 & 9,782 (43)   & $-$           & 10,582 (109) \\
Nissum Fjord, Felsted Kog  & 196 & 9,391  & 9,232 (49)   & 8,698 (38)   & 9,666 (109) \\
Thisted Bredning           & 244 & 9,605  & 9,858 (50)   & 9,692 (27)   & 9,527 (167) \\
Nibe Bredning og Langerak  & 704 & 8,849  & 8,803 (229)  & 8,580 (75)   & 8,917 (400) \\
Randers Fjord, indre       & 807 & 9,296  & 9,121 (200)  & 9,256 (88)   & 9,363 (519) \\
Horsens Fjord, indre       & 233 & 9,638  & 9,431 (57)   & 9,895 (32)   & 9,666 (144) \\
Odense Fjord, Seden Strand & 304 & 9,057  & 8,998 (73)   & 8,866 (40)   & 9,104 (191) \\
Isefjord, indre            & 202 & 9,331  & 8,937 (45)   & $-$           & 9,486 (139) \\
Karrebæk Fjord             & 342 & 8,825  & 9,028 (71)   & 8,698 (50)   & 8,789 (221) \\
Roskilde Fjord, indre      &  90 & 9,215  & 9,339 (23)   & $-$           & 9,190 (57) \\
\addlinespace\bottomrule
\end{tabular}
\begin{tablenotes}
\footnotesize
\item \textit{Notes:} All values are measured in DKK per hectare, and weighted by area;
Costs denotes the catchment-level area-weighted average total costs (Costs I + Costs II);
Fine, Coarse, and Loamy refer to average total costs for fields on the respective
soil type within each catchment; numbers in parentheses indicate the number of
fields; `$-$' indicates fewer than 20 fields and is therefore not reported.
\end{tablenotes}
\end{threeparttable}
\end{table}



Tabel~\ref{tab:costs-vinterhvede} show average total costs (DKK/ha) for winter wheat fields by catchment and soil type. Overall, total production costs exhibit substantially less spatial variation than gross yields, and no clear west–east gradient is observed. Instead, costs are relatively stable across catchments, typically ranging between 8,600 and 9,666 DKK/ha. A notable exception is Nissum Bredning, which exhibits consistently higher costs across soil types (9,782---10,582 DKK/ha), indicating persistently higher production expenses relative to other catchments. In contrast, Karrebæk Fjord combines relatively low costs with high yields, suggesting comparatively favourable production conditions.


\begin{table}[H]
\centering
\small
\caption{Average profitability (DKK/ha/year) for winter wheat fields by catchment and soil type}
\label{tab:dbii-vinterhvede}
\begin{threeparttable}
\begin{tabular}{lrrrrr}
\toprule
\textbf{Catchment} & \textbf{Fields} & \textbf{DBII} & \textbf{Fine} &
\textbf{Coarse} & \textbf{Loamy} \\
\midrule\addlinespace
Lister Dyb                 & 440 &  $-$74   & $-$501 (117)  &    22 (71)    &    89 (252) \\
Nissum Bredning            & 169 & $-$2,071 & $-$657 (43)   & $-$            & $-$2,669 (109) \\
Nissum Fjord, Felsted Kog  & 196 & $-$1,238 & $-$1,126 (49) & $-$912 (38)   & $-$1,383 (109) \\
Thisted Bredning           & 244 &   $-$4   & $-$530 (50)   & $-$1,155 (27) &    289 (167) \\
Nibe Bredning og Langerak  & 704 &   621    &   557 (229)   &   167 (75)    &    729 (400) \\
Randers Fjord, indre       & 807 &   783    &   976 (200)   &   529 (88)    &    749 (519) \\
Horsens Fjord, indre       & 233 &    86    &   240 (57)    & $-$634 (32)   &    150 (144) \\
Odense Fjord, Seden Strand & 304 & 1,637    & 1,591 (73)    &   991 (40)    & 1,743 (191) \\
Isefjord, indre            & 202 & 1,597    & 2,692 (45)    & $-$            & 1,310 (139) \\
Karrebæk Fjord             & 342 & 2,274    & 2,353 (71)    & 1,567 (50)    & 2,362 (221) \\
Roskilde Fjord, indre      &  90 & 1,304    & 1,031 (23)    & $-$            & 1,569 (57) \\
\addlinespace\bottomrule
\end{tabular}
\begin{tablenotes}
\footnotesize
\item \textit{Notes:} All values are measured in DKK per hectare, and weighted by area;
DBII denotes gross margin II; Fine, Coarse, and Loamy refer to DBII for fields on the
respective soil type within each catchment; numbers in parentheses indicate the number of
fields; `$-$' indicates fewer than 20 fields and is therefore not reported.
\end{tablenotes}
\end{threeparttable}
\end{table}


Table~\ref{tab:dbii-vinterhvede} presents area-weighted average profitability (DBII) for winter wheat across catchments and soil types.  The results show substantial spatial variation in profitability, with both positive and negative average DBII observed across catchments. Seven catchments exhibit positive average DBII: Karrebæk Fjord (+2,274 DKK/ha), Odense Fjord (+1,637 DKK/ha), Isefjord (+1,597 DKK/ha), Roskilde Fjord (+1,304 DKK/ha), Randers Fjord (+783 DKK/ha), Nibe Bredning og Langerak (+621 DKK/ha), and Horsens Fjord (+86 DKK/ha). The remaining four catchments are negative, with Nissum Bredning ($-$2,071 DKK/ha) and Nissum Fjord ($-$1,238 DKK/ha) the most negative. Overall, there is considerable heterogeneity across catchments, with a difference of more than 4,300 DKK/ha across all catchments. Excluding Nissum Bredning og Langerak and Karrebæk Fjord, which represent extreme values at each end of the distribution, the range is reduced to 2,875 DKK/ha. This indicates considerable spatial heterogeneity in winter wheat profitability.


The west--east gradient observed in yields is also reflected in profitability, consistent with relatively limited spatial variation in costs. Soil type variation broadly follows expected agronomic patterns --- loamy soils generally outperform coarse sand --- however, within-catchment variation is smaller than variation across catchments. Two exceptions are worth noting. In Nissum Bredning, fine sand ($-$657 DKK/ha) outperforms loamy soil ($-$2,669 DKK/ha). In Isefjord, fine sand (+2,692 DKK/ha) substantially outperforms loamy soil (+1,310 DKK/ha), though this estimate rests on only 45 observations and should be interpreted with caution.

Overall, spatial variation in winter wheat profitability seems primarily driven by differences in yields. Catchments in western Jutland combine relatively low yields with moderate to high costs, resulting in negative DBII, whereas eastern Jutland and Zealand exhibit higher yields combined with comparable or lower costs, leading to substantially higher profitability. Variation in costs plays a secondary role and does not follow a clear geographic pattern. The results are consistent with the expected agronomic gradient between eastern and western Denmark. However, they may also reflect weather conditions in 2024, in particular excessive rainfall in western and southern Denmark (see Figure~\ref{fig:dmi_nedbor}). For this reason, the results should be interpreted as specific to 2024. 





\subsubsection{Field-Level Profitability for Spring Barley in 2024} 


\begin{table}[H]
\centering
\small
\caption{Average spring barley gross yield (DKK/ha/year) by catchment and soil type}
\label{tab:yield-vaarbyg}
\begin{threeparttable}
\begin{tabular}{lrrrrr}
\toprule
\textbf{Catchment} & \textbf{Fields} & \textbf{Yield} &
\textbf{Fine} & \textbf{Coarse} & \textbf{Loamy} \\
\midrule\addlinespace
Lister Dyb                 & 718  & 5,871 & 5,916 (170) & 5,600 (114) & 5,930 (434) \\
Nissum Bredning            & 434  & 6,517 & 6,510 (104) & 6,382 (67)  & 6,545 (263) \\
Nissum Fjord, Felsted Kog  & 910  & 6,804 & 6,930 (247) & 6,274 (165) & 6,912 (498) \\
Thisted Bredning           & 329  & 6,125 & 5,525 (96)  & 5,858 (49)  & 6,418 (184) \\
Nibe Bredning og Langerak  & 1,031& 6,418 & 6,779 (362) & 5,506 (149) & 6,380 (520) \\
Randers Fjord, indre       & 926  & 6,462 & 6,408 (222) & 6,080 (125) & 6,554 (579) \\
Horsens Fjord, indre       & 177  & 5,906 & 6,116 (36)  & 5,749 (23)  & 5,872 (118) \\
Odense Fjord, Seden Strand & 247  & 6,525 & 6,834 (58)  & 5,425 (36)  & 6,729 (153) \\
Isefjord, indre            & 109  & 6,377 & 5,796 (25)  & $-$          & 6,638 (75) \\
Karrebæk Fjord             & 256  & 5,710 & 5,907 (58)  & 4,915 (32)  & 5,763 (166) \\
Roskilde Fjord, indre      & 59   & 4,647 & $-$          & $-$          & 4,916 (37) \\
\addlinespace\bottomrule
\end{tabular}
\begin{tablenotes}
\footnotesize
\item \textit{Notes:} Values are measured in kg per hectare and weighted by area;
Fine, Coarse, and Loamy denote average yields for fields on the respective soil type;
numbers in parentheses indicate the number of fields; `$-$' indicates fewer than 20 fields and is therefore not reported.
\end{tablenotes}
\end{threeparttable}
\end{table}



Table~\ref{tab:yield-vaarbyg} presents average spring barley yields across 11 catchments in 2024. In contrast to winter wheat, which exhibits a clear west–east productivity gradient, spring barley yields shows no consistent geographic pattern. Excluding Roskilde Fjord, indre, yields are relatively similar across catchments 1,094 DKK/ha —-- from 5,710 DKK/ha (Karrebæk Fjord) to 6,804 DKK/ha (Nissum Fjord, Felsted Kog) --- and no systematic spatial structure is observed.

Roskilde Fjord, indre stands out, with a substantially lower average value (4,647 DKK/ha) based on 59 fields. Given the limited sample size, this estimate should be interpreted with caution.

Within catchments, coarse soils consistently exhibit the lowest yields, while fine and loamy soils perform comparably. The gap is particularly pronounced in Nibe Bredning og Langerak, where coarse soils yield 5,506 DKK/ha against 6,779 DKK/ha on fine soils, and in Odense Fjord, Seden Strand, where the corresponding figures are 5,425 DKK/ha and 6,834 DKK/ha. These within-catchment differences suggest that soil texture is a meaningful driver of spring barley productivity. 


Comparing winter wheat and spring barley, within-catchment yield differences across soil types are of broadly similar magnitude for the two crops. 
Comparing catchments where all three soil types are observed for both crops, the average difference between the highest- and lowest-yielding soil type is approximately 799 DKK/ha for spring barley and 846 DKK/ha for winter wheat. Neither crop consistently exhibits greater soil-type sensitivity than the other. In some catchments, such as Nibe Bredning og Langerak and Odense Fjord, Seden Strand, the difference is more pronounced for spring barley, while in others, such as Lister Dyb and Thisted Bredning, winter wheat shows greater within-catchment variation. 

Across both crops, coarse soils consistently yield the lowest values within catchments, while the relative performance of fine and loamy soils is more variable. However, this pattern is likely partly reflecting the structure of the allocation model used to assign yields.



\begin{table}[H]
\centering
\small
\caption{Average costs (DKK/ha/year) for spring barley fields by catchment and soil type}
\label{tab:costs-vaarbyg}
\begin{threeparttable}
\begin{tabular}{lrrrrr}
\toprule
\textbf{Catchment} & \textbf{Fields} & \textbf{Costs} &
\textbf{Fine} & \textbf{Coarse} & \textbf{Loamy} \\
\midrule\addlinespace
Lister Dyb                 & 718 & 8,483 & 8,418 (170) & 8,612 (114) & 8,471 (434) \\
Nissum Bredning            & 434 & 9,514 & 9,204 (104) & 9,781 (67)  & 9,599 (263) \\
Nissum Fjord, Felsted Kog  & 910 & 8,702 & 8,702 (247) & 8,703 (165) & 8,701 (498) \\
Thisted Bredning           & 329 & 8,333 & 8,274 (96)  & 8,296 (49)  & 8,363 (184) \\
Nibe Bredning og Langerak  & 1,031 & 8,386 & 8,342 (362) & 7,844 (149) & 8,542 (520) \\
Randers Fjord, indre       & 926 & 8,609 & 8,625 (222) & 8,795 (125) & 8,569 (579) \\
Horsens Fjord, indre       & 177 & 9,304 & 9,071 (36)  & 9,248 (23)  & 9,394 (118) \\
Odense Fjord, Seden Strand & 247 & 8,925 & 9,232 (58)  & 8,007 (36)  & 9,075 (153) \\
Isefjord, indre            & 109 & 9,419 & 8,312 (25)  & $-$          & 9,757 (75) \\
Karrebæk Fjord             & 256 & 8,432 & 8,617 (58)  & 8,184 (32)  & 8,412 (166) \\
Roskilde Fjord, indre      & 59  & 8,875 & $-$          & $-$          & 8,937 (37) \\
\addlinespace\bottomrule
\end{tabular}
\begin{tablenotes}
\footnotesize
\item \textit{Notes:} Values are measured in DKK per hectare and weighted by area;
Fine, Coarse, and Loamy denote average costs for fields on the respective soil type;
numbers in parentheses indicate the number of fields; `$-$' indicates fewer than 20 fields and is therefore not reported.
\end{tablenotes}
\end{threeparttable}
\end{table}


Table~\ref{tab:costs-vaarbyg} presents average production costs for spring barley across catchments. Catchment-level costs range from 8,333 DKK/ha (Thisted Bredning) to 9,514 DKK/ha (Nissum Bredning), corresponding to a difference of 1,181 DKK/ha. As with yields, costs show no clear geographic pattern across catchments.


Within-catchment variation accross soil types is generally modest. Nissum Fjord, Felsted Kog is a notable case: where costs are virtually identical across all three soil types (8,701 -8,703), making it difficult to identify any soil-related differences. More broadly, there is no cosistent directional pattern in costs across soil types --- coarse soils are the most expensive in some catchments, such as Nissum Bredning (9,781 DKK/ha) and Randers Fjord (8,795 DKK/ha), while they are among the least expensive in others, such as Nibe Bredning og Langerak (7,844 DKK/ha) and Odense Fjord, Seden Strand (8,007 DKK/ha). 

This contrasts with yields, where coarse soils consistently exhibit the lowest returns within catchments. Overall, these patterns mirror the findings for winter wheat, suggesting that soil type does not exert a consistent influence on the cost structure.


\begin{table}[H]
\centering
\small
\caption{Average profitability (DBII, DKK/ha/year) for spring barley fields by catchment and soil type}
\label{tab:dbii-vaarbyg}
\begin{threeparttable}
\begin{tabular}{lrrrrr}
\toprule
\textbf{Catchment} & \textbf{Fields} & \textbf{DBII} &
\textbf{Fine} & \textbf{Coarse} & \textbf{Loamy} \\
\midrule\addlinespace
Lister Dyb                 & 718 & $-$2,612 & $-$2,502 (170) & $-$3,012 (114) & $-$2,541 (434) \\
Nissum Bredning            & 434 & $-$2,997 & $-$2,694 (104) & $-$3,399 (67)  & $-$3,055 (263) \\
Nissum Fjord, Felsted Kog  & 910 & $-$1,898 & $-$1,772 (247) & $-$2,429 (165) & $-$1,789 (498) \\
Thisted Bredning           & 329 & $-$2,207 & $-$2,749 (96)  & $-$2,438 (49)  & $-$1,945 (184) \\
Nibe Bredning og Langerak  & 1,031 & $-$1,967 & $-$1,563 (362) & $-$2,339 (149) & $-$2,163 (520) \\
Randers Fjord, indre       & 926 & $-$2,147 & $-$2,217 (222) & $-$2,715 (125) & $-$2,015 (579) \\
Horsens Fjord, indre       & 177 & $-$3,398 & $-$2,955 (36)  & $-$3,499 (23)  & $-$3,522 (118) \\
Odense Fjord, Seden Strand & 247 & $-$2,399 & $-$2,398 (58)  & $-$2,582 (36)  & $-$2,347 (153) \\
Isefjord, indre            & 109 & $-$3,042 & $-$2,517 (25)  & $-$             & $-$3,119 (75) \\
Karrebæk Fjord             & 256 & $-$2,722 & $-$2,710 (58)  & $-$3,269 (32)  & $-$2,649 (166) \\
Roskilde Fjord, indre      & 59  & $-$4,228 & $-$             & $-$             & $-$4,021 (37) \\
\addlinespace\bottomrule
\end{tabular}
\begin{tablenotes}
\footnotesize
\item \textit{Notes:} All values are measured in DKK per hectare and weighted by area;
DBII denotes gross margin II; Fine, Coarse, and Loamy refer to DBII for fields on the
respective soil type; numbers in parentheses indicate the number of fields; `$-$' indicates fewer than 20 fields and is therefore not reported.
\end{tablenotes}
\end{threeparttable}
\end{table}

Table~\ref{tab:dbii-vaarbyg} presents average profitability (DBII) from spring barley across catchments. The most striking result is that all catchments exhibit negative average DBII in 2024. 

Excluding Roskilde Fjord, indre, due to the limited observations, DBII ranges from −1,898 DKK/ha (Nissum Fjord, Felsted Kog) to −3,398 DKK/ha (Horsens Fjord, indre), corresponding to a spread of approximately 1,500 DKK/ha. As with yields and costs, no clear geographic pattern in profitability is observed across catchments.

Roskilde Fjord, indre represents a clear outlier, with a DBII of −4,228 DKK/ha based on only 59 fields. This estimate should therefore be interpreted with caution.

Within catchments, coarse soils consistently exhibit the lowest profitability in all catchments with full soil-type coverage. This pattern reflects the lower yields of coarse soils documented above.

Overall, the spring barley results are more difficult to interpret than those for winter wheat. Most notably, all catchments record negative DBII in 2024, which may partly reflect adverse weather conditions during the growing season (Section~\ref{sec:Weather Conditions}), where spring barley yields were approximately 9.1 percent below the average (Section~\ref{sec:SelectedCrops}). 

 

It should be noted that the initial results were based on crop-producing farms only. In that subsample, a spatial gradient was  observed, indicating higher DKK/ha values in livestock-dense regions in Jutland. When the full sample is included, this pattern disappears. This suggests that spring barley values may be systematically higher in crop-specialised farms compared to mixed farms in these catchments (refer to appendix xx).

One possible explanation is differences in output valuation. Spring barley is predominantly used as on-farm animal feed, and in livestock-dense catchments a larger share of production may be recorded at internal transfer prices rather than market prices. If internal transfer prices are systematically lower than market prices, this would generate a downward bias in observed yields in livestock-intensive catchments.

Although this hypothesis cannot be tested with the available data, it suggests that variation in internal pricing of feed crops may be an important source of measurement bias. The potential divergence between market prices and internally valued feed crops therefore warrants further investigation.  



\end{document}